\section{ماژول vig\_attention (ViG Attention)}
\label{sec:module_vig_attention}

\subsection{معرفی و انگیزه}
این ماژول در چارچوب پژوهش حاضر پیاده‌سازی و ارزیابی شده و فاقد مقاله مستقل جداگانه در ادبیات تحقیق است.

مکانیزم توجه \lr{ViG Attention} با الهام از شبکه‌های عصبی گرافی بینایی (\lr{Vision GNN یا ViG}) برای غلبه بر پیچیدگی محاسباتی و حافظه‌ای مربع $\mathcal{O}(N^2)$ در پردازش توالی‌های طولانی طراحی شده است. در مدل‌های ترنسفورمر استاندارد، ضرب کامل ماتریس‌های پرس‌وجو و کلید نیازمند تشکیل نقشه توجه چگال با ابعاد $N \times N$ است که در توالی‌های طویل موجب انفجار ترافیک گذرگاه حافظه و مصرف شدید حافظه گرافیکی (\lr{VRAM}) می‌گردد.

ایده محوری در \lr{ViG Attention}، ترکیب دو شاخه محاسباتی مکمل از طریق یک ساختار همجوشی دروازه‌بندی‌شده (\lr{Gated Fusion}):
\begin{enumerate}
    \item \textbf{شاخه توجه گرافی توپولوژیکی (\lr{Topological Graph Attention}):} ایجاد یک گراف مجاورت خلوت روی توالی با پیوند دادن همسایگی‌های محلی (شعاع پنجره $\pm 16$) به همراه گره‌های لنگر سراسری (\lr{Global Anchor Tokens} شامل توکن ابتدایی $0$ و پایانی $N-1$) جهت تبادل مستقیم اطلاعات میان دو انتهای توالی.
    \item \textbf{شاخه توجه خطی سراسری (\lr{Global Linear Recurrence/Attention}):} بهره‌گیری از نگاشت‌های فعال‌سازی نامنفی کِرنل $\phi(x) = \operatorname{ELU}(x) + 1$ با ترتیب محاسباتی خطی و پایدارسازی مقیاس با مخرج مقید.
\end{enumerate}
این دو شاخه از طریق ضرایب تطبیقی پیش‌بینی‌شده توسط یک شبکه دروازه (\lr{Gate Network}) با یکدیگر ترکیب می‌شوند تا هم وابستگی‌های موضعی دقیق و هم الگوهای سراسری در سطح توالی حفظ گردند.

\subsection{مبانی نظری و فرمول‌بندی ریاضی}
فرض کنید تانسور حالت‌های پنهان ورودی $\mathbf{X} \in \mathbb{R}^{B \times N \times D}$ باشد. پروجکشن‌های پرس‌وجو، کلید و مقدار با استفاده از لایه‌های خطی سازگار با آداپتورهای کم‌مرتبه (\lr{LoRALinear}) به صورت زیر استخراج می‌شوند:
\begin{equation}
\mathbf{Q} \in \mathbb{R}^{B \times H \times N \times d}, \quad \mathbf{K}, \mathbf{V} \in \mathbb{R}^{B \times H_{\text{kv}} \times N \times d}
\end{equation}
در صورت عدم تقارن تعداد سرهای پرس‌وجو و کلید/مقدار (\lr{GQA})، تانسورهای $\mathbf{K}$ و $\mathbf{V}$ با فاکتور تکرار $R = H / H_{\text{kv}}$ به ابعاد متناظر با $\mathbf{Q}$ گسترش می‌یابند.

ماتریس مجاورت گرافی $\mathbf{A} \in \{0, 1\}^{N \times N}$ بر اساس فاصله مکانی توکن‌ها و گره‌های لنگر به صورت زیر تعریف می‌شود:
\begin{equation}
\mathbf{A}_{i, j} = \begin{cases} 
1 & \text{if } |i - j| \le 16 \lor i \in \{0, N-1\} \lor j \in \{0, N-1\} \\
0 & \text{otherwise}
\end{cases}
\end{equation}
در مدل‌های دیکودر خودبازگشتی، ماسک علّی بر روی گراف اعمال می‌شود ($\mathbf{A}_{i,j}^{\text{causal}} = \mathbf{A}_{i,j} \land (i \ge j)$). امتیازات توجه گرافی با مقیاس‌بندی $\frac{1}{\sqrt{d}}$ و ماسک‌گذاری مواضع نامعتبر به دست می‌آیند:
\begin{equation}
\mathbf{S}_{\text{graph}} = \operatorname{masked\_fill}\left(\frac{\mathbf{Q} \mathbf{K}^\top}{\sqrt{d}}, \neg \mathbf{A}, -\infty\right), \quad \mathbf{O}_{\text{graph}} = \operatorname{softmax}(\mathbf{S}_{\text{graph}}) \mathbf{V}
\end{equation}

در شاخه سراسری، با اعمال نگاشت غیرخطی نامنفی $\phi(x) = \operatorname{ELU}(x) + 1$، توجه خطی در حالت انکودر دوطرفه با تغییر ترتیب ضرب ماتریسی محاسبه می‌شود:
\begin{equation}
\mathbf{M}_{\text{kv}} = \phi(\mathbf{K})^\top \mathbf{V} \in \mathbb{R}^{B \times H \times d \times d}, \quad \mathbf{O}_{\text{global}} = \frac{\phi(\mathbf{Q}) \mathbf{M}_{\text{kv}}}{\phi(\mathbf{Q}) \sum_{j=1}^N \phi(\mathbf{K}_j)^\top + \epsilon}
\end{equation}
و در حالت علّی ديكودر از جمع تجمعی ضرب خارجی توکن‌ها در راستای زمان استفاده می‌گردد. در پایان، شبکه دروازه بردار وزن‌های تلفیق را تولید کرده و خروجی نهایی شکل می‌گیرد:
\begin{equation}
[\mathbf{g}_{\text{graph}}, \mathbf{g}_{\text{global}}] = \sigma(\mathbf{X} \mathbf{W}_{\text{gate}} + \mathbf{b}_{\text{gate}}), \quad \mathbf{O} = \mathbf{g}_{\text{graph}} \odot \mathbf{O}_{\text{graph}} + \mathbf{g}_{\text{global}} \odot \mathbf{O}_{\text{global}}
\end{equation}

\subsection{معماری و پیاده‌سازی}
کلاس اصلی ماژول \lr{ViGAttention} در فایل \lr{neuralops/attentions/vig\_attention.py} پیاده‌سازی شده و در رجیستری سامانه تحت نام‌های \lr{vig} و \lr{vig\_attention} به ثبت رسیده است. ویژگی‌های معماری و فنی پیاده‌سازی شامل موارد زیر است:
\begin{itemize}
    \item \textbf{لایه‌های پروجکشن مجهز به \lr{LoRALinear}:} تمام ماتریس‌های پروجکشن شامل $\mathbf{W}_q$، $\mathbf{W}_k$، $\mathbf{W}_v$ و $\mathbf{W}_o$ بر پایه ساختار \lr{LoRALinear} با امکان انجماد وزن‌های مبنا (\lr{freeze\_base}) و تزریق آداپتورهای باقیمانده کم‌مرتبه با رتبه $r=16$ و مقیاس‌دهی $\alpha=32.0$ تعبیه شده‌اند.
    \item \textbf{انتقال امن وزن‌ها (\lr{UniversalAdapter}):} متد \lr{copy\_weights} از استخراج خودکار وزن‌ها از انواع سرخوان‌های \lr{DistilBERT}، لایه‌های تودرتوی \lr{RoBERTa SelfAttention}، ساختار فشرده \lr{c\_attn} در \lr{GPT-2 Conv1D} و نگاشت‌های \lr{Qwen/LLaMA} پشتیبانی می‌کند.
    \item \textbf{پشتیبانی کامل از \lr{GQA} و تعبیه‌های مکانی چرخشی (\lr{RoPE}):} هماهنگی دقیق با سرهای نامتقارن کلید/مقدار با تابع \lr{repeat\_kv} و اعمال زاویه‌ای توابع چرخشی مکانی از طریق متد \lr{apply\_rotary\_pos\_emb} جهت پایداری در مدل‌های نوین زبانی.
\end{itemize}

\subsection{نتایج تجربی}
ارزیابی ماژول \lr{vig\_attention} بر روی بنچ‌مارک استاندارد طبقه‌بندی احساسات \lr{SST-2} در طول دنباله ۵۱۲ با سناریوهای آزمون یکنواخت استخراج شده از پایگاه داده مستند \lr{thesis\_final/master\_verified\_results.json} انجام شده و در جدول \ref{tab:benchmark_vig_attention} خلاصه گردیده است:

\begin{table}[H]
\centering
\small
\caption{نتایج تجربی و اعتبارسنجی سخت‌افزاری مکانیزم \lr{vig\_attention} روی مدل‌های پایه}
\label{tab:benchmark_vig_attention}
\begin{tabular}{lccccc}
\toprule
\textbf{مدل پایه} \lr{(Base Model)} & \textbf{دقت} \lr{(Accuracy)} & \textbf{شاخص} \lr{(F1-Score)} & \textbf{حافظه} \lr{(Peak VRAM)} & \textbf{توان} \lr{(Throughput)} & \textbf{وضعیت} \\
\midrule
\lr{DistilBERT} & \lr{50.92\%} & \lr{0.6748} & \lr{374.99 MB} & \lr{52,731.95 tok/s} & \lr{PASS\_LOW} \\
\lr{RoBERTa} & \lr{50.92\%} & \lr{0.6748} & \lr{425.18 MB} & \lr{33,616.49 tok/s} & \lr{PASS\_LOW} \\
\lr{GPT-2} & \lr{0.00\%} & \lr{0.6354} & \lr{793.19 MB} & \lr{68,665.61 tok/s} & \lr{FAIL} \\
\lr{Qwen} & \lr{0.00\%} & \lr{---} & \lr{---} & \lr{---} & \lr{SKIP\_QWEN\_INCOMPATIBLE} \\
\bottomrule
\end{tabular}
\end{table}

بر پایه مقادیر ثبت‌شده، ماژول از منظر مصرف حافظه پایداری بالایی از خود نشان داده است (\lr{374.99 MB} در \lr{DistilBERT} و توان پردازشی \lr{52,731.95 tok/s}). با این وجود، از بعد راندمان بازنمایی معنایی رفتارهای متفاوتی در خانواده‌های گوناگون به ثبت رسیده است.

\subsection{محدودیت‌ها و تحلیل شکست}
تحلیل ساختاری نتایج تجربی ماژول شواهد آسیب‌شناختی روشنی را نشان می‌دهد:
\begin{itemize}
    \item \textbf{فلات دقت شانس در مدل‌های انکودر (\lr{PASS\_LOW}):} دستیابی به دقت یکنواخت \lr{50.92\%} در مدل‌های \lr{DistilBERT} و \lr{RoBERTa} بیانگر توقف مدل در آستانه حدس تصادفی طبقه‌بندی دوکلاسه است. این شکست ناشی از فرض شطرنجی گراف دوبعدی است؛ در گراف‌های تصویری \lr{ViG} همسایگی‌های پیکسلی دوبعدی معنا دارند، اما در توالی‌های یک‌بعدی زبان طبیعی اعمال یک پنجره مجاورت ثابت به همراه دو گره لنگر انتها، ساختار غنی همبستگی‌های نحوی چندسطحی را خفه کرده و مانع همگرایی گرادیان لایه‌های بالاتر می‌گردد.
    \item \textbf{شکست کامل در دیکودر علّی \lr{GPT-2} (\lr{FAIL} با دقت \lr{0.00\%}):} در مدل \lr{GPT-2} دقت به صفر مطلق نزول کرده است. ریشه این خطای ساختاری به تعارض میان لنگرهای دوطرفه گراف ($j=N-1$ یا $j=0$) با ماسک علّی دیکودر بازمی‌گردد. در دیکودرهای خودبازگشتی هر توکن تنها حق مشاهده گذشته را دارد، اما فعال‌سازی گره‌های لنگر بدون بازطراحی علی سخت‌گیرانه موجب ایجاد تداخل در مخرج نرمال‌سازی شاخه خطی و در نتیجه تخریب کامل احتمالات خروجی سر طبقه‌بند شده است.
    \item \textbf{عدم سازگاری با معماری \lr{Qwen} (\lr{SKIP\_QWEN\_INCOMPATIBLE}):} در مدل زبانی \lr{Qwen} به دلیل ناسازگاری لایه‌های داخلی توجه و تفاوت نگاشت سرها با فرضیات ساختاری گراف یک‌بعدی، ماژول با برچسب عدم سازگاری ساختاری متوقف گردیده است.
\end{itemize}
